\section{Data and Sample Construction}\label{sec:data}

The data come from BEC, the state of S\~ao Paulo's centralized electronic procurement platform. BEC intermediates all standardized-goods procurement for state bureaus and affiliated units, so the raw extract covers the universe of item-level bids, transacted prices, firm identifiers, quantities, and auction identifiers for January 2016 through December 2019---3.7~million observations across 82{,}569 items, 832{,}984 purchase orders, and 1{,}344 public buyer units. Group~65 (medical, dental, and hospital supplies) accounts for roughly 27 percent of transactions over the window and is the treated group; the remaining 76 product groups form the never-treated control set in the difference-in-differences benchmark of Appendix~\ref{app:did}.

The structural analysis narrows to a sharper sample. Identification uses a symmetric window of 18 months on each side of the March~2018 cutoff (36 months in total, the ``18m'' label of Table~\ref{tab:v3_window_sensitivity}); in Stata monthly-date units the window is $[680, 715]$, or September~2016 through August~2019 inclusive. Within that window, I restrict to Preg\~ao auctions in Group~65. The restriction does two things. Preg\~ao is the descending-clock format that admits point identification of cost distributions from drop-out bids, so the structural estimation requires it; Group~65, in turn, is where the policy bound---the only product group to switch regime during the sample.

Bids are organized at the firm-auction level. Each observation is the final bid submitted by one firm in one item-auction, together with a winner/loser flag. I apply three filters. First, bids must satisfy $c_\epsilon = b / p^{\mathrm{ref}} \in (0, 3]$; the upper cutoff removes implausible outliers (less than 0.4 percent of bids, almost always data-entry errors at the hundred-times-reference level). Second, auctions must have at least two participating firms---below that, neither point identification of $F_c$ nor the BNE counterfactual is well-defined. Third, each bid carries a type flag (SME or non-SME, from the BEC administrative SME registry, validated against historical bidding to avoid misclassifying firms that temporarily exited SME status) and a pharma flag (CADMAT classes 6531, 6532, 6536, 6581---CMED-regulated pharmaceuticals and biologicals---are pharma-narrow; the remainder of Group~65 is non-pharma).

The resulting structural sample is 297{,}967 firm-auction observations across 97{,}993 distinct auctions, balanced across Pre (September~2016 to February~2018, 18 months, $n = 48{,}740$ auctions) and Post (March~2018 to August~2019, 18 months, $n = 49{,}253$). Descriptive counts by stratum appear in Appendix Tables~\ref{tab:v3_s1_handoff} and~\ref{tab:v3_pharma_counts}.

Auction-level unobserved heterogeneity is material and not fully absorbed by the normalization against the item-level reference price. Two items inside the same CADMAT class---say two lots of the same dressing size differing in volume, or two pharma items with ANVISA registration costs of very different orders of magnitude---share a scale component that moves every bid in the auction together. I handle this with a method-of-moments variance decomposition followed by best-linear-predictor shrinkage of the auction-level component out of each log-bid, following \citet{krasnokutskaya2011}. The shrinkage recovers intraclass correlations of $\rho^k$ that average roughly 0.28--0.59 across the four Preg\~ao (pharma $\times$ type) cells---the structural sample---and 0.55--0.66 across the corresponding Convite cells. The lower Preg\~ao bound is pulled down by the pharmaceutical-SME-Pre stratum ($n = 10{,}791$ bids), where the auction-level variance $\sigma_a^2$ is small relative to within-auction noise and the resulting cell-specific $\rho$ is close to zero; restricted to the three remaining Preg\~ao cells the range is 0.36--0.59, comparable to the 0.50--0.70 range Krasnokutskaya (2011) reports for California highway procurement. The cleaned residuals $c_\epsilon^{\text{clean}} = \exp(\hat e_{it})$, where $\hat e_{it}$ is the bidder-specific log-bid residual after the auction-level component $\hat a_t$ is shrunk out, enter all downstream structural estimation; the corresponding variance panel is Appendix Table~\ref{tab:v3_uh_variance}. Section~\ref{sec:identification} develops the statistical details.

Entry counts enter separately. For each auction I count the distinct firms of each type that submitted at least one bid, then average within (period $\times$ pharma) cells. In Preg\~ao, non-pharma moves from 0.94 SMEs and 2.68 non-SMEs per auction pre-policy to 1.87 SMEs and 1.50 non-SMEs post-policy. Pharma moves from 0.55 and 2.61 to 1.22 and 1.66. The roughly two-fold SME increase is the object of the extensive-margin channel in the counterfactual; the accompanying non-SME decline is the composition shift the set-aside induces. Entry rates enter the BNE as observed equilibrium arrivals (Section~\ref{sec:identification}); the panel is Appendix Table~\ref{tab:v3_entry_rates}.

Four ancillary inputs anchor the analysis. Reference prices come from the BEC catalog (\texttt{preco\_ref}) and are used both for normalization and for pricing out the fiscal cost of the rule in reais. The reference price for each item is set ex ante by the buyer from the BEC catalog and does not move with the regime change: item-level fixed effects in the DiD specifications and the structural normalization absorb cross-sectional variation in $p^{\mathrm{ref}}$ across items, and the BEC catalog architecture sets $p^{\mathrm{ref}}$ at the item level for the entire window. The variable that changes around the cutoff is regime admissibility, not item-level catalog pricing. CADMAT class codes come from the BEC product hierarchy. The historical SME flag is reconstructed from firm-level bidding histories over the full 2014--2019 period to protect against transient status changes. The pharma flag follows the CADMAT class list above.

The pair of furosemide 40~mg purchases mentioned in the introduction (item 110639, a CMED-regulated generic antihypertensive) sits at the institutional level. In February~2018, one month before the cutoff, a S\~ao Paulo public buyer procured fifty thousand units under the open regime: ten firms bid, and the clearing price settled roughly twelve percent below the reference. In October~2018, eight months after the cutoff, a comparable purchase by the same buyer attracted three SMEs, and the winning price cleared just above reference. The pair is illustrative; the paper's pharmaceutical estimates are averages over 59{,}997 Preg\~ao pharmaceutical auctions in the 18-month structural sample (Table~\ref{tab:v3_s1_handoff}).
